\subsection{Research findings}
\label{sec:findings}

The findings from the experiments described in \hyperref[sec:simulation_platform]{Section 3.3.1} are discussed in the context of the research objectives and what
insights can be gleaned by framing the results of testing the FLISR Edge solution in terms of these objectives.

\begin{itemize}
  \item \textbf{R1 - How effective will the outputs of the FLISR solution be in terms of reducing DSO fault KPIs?}
\end{itemize}

\noindent From simulating a variety of fault cases utilising the FLISR Edge simulation platform, it has been found that the reduction of DSO fault KPIs, such as CML, varies on a case by
case basis. For example, test cases where there are a low vs high occurrence of faulted switch devices, the difference in affected customers and the resulting CML generated per hour can be drastic.
Additionally, results for networks which are less complex and options for back fed supply are limited, such as Kerry for example, the results are less impressive.
This observation supports the research literature where it is found that FLISR is most effective on complex, radial networks, such as the network located in Waterford \cite{7000578}.
\\
In addition, it is important to note that FLISR is a fault mitigation solution, therefore in cases where a large number of faulted lines occur on a section of network, the results will be
limited. For example, a specific test case simulated for the Waterford trial where one switch is set to faulted, results in 1370 affected customers and 78090 CML per hour prior to FLISR operation,
post operation this is reduced to 171 customers affected and 9747 CML per hour. However, increasing the number of faulted switches to four for this same test case,
results in an affected customer count of 844 after operation and a total of 48108 CML after FLISR operation.

\begin{itemize}
  \item \textbf{R2 - Given a similar set of conditions, how will the results produced by the centralised FLISR algorithm compare to the distributed FLISR algorithms operating locally at the edge of the network?}
\end{itemize}

\noindent To determine the answer to this question, a wide variety of test cases were simulated across the three distribution grids in centralised and distributed FLISR modes and these were compared and contrasted.
From the results gathered, the operational steps generated by the algorithms are identical, with no observable difference in terms of the steps taken to resolve a given fault event.
In practice, this is a desirable outcome, as when applied to the field, divergent control operations generated by each algorithm may suggest that the algorithms require refactoring.
Furthermore, it is not definitive at this stage that identical results will be given for every situation, as only three networks have been tested in a simulated environment, in addition,
with further testing and validation other factors and distribution network features may be found which impact the FLISR algorithm logic.
\\
One key difference observed between the centralised and distributed FLISR algorithms is the operational time. As the centralised algorithm requires a network connection to ingest and output generated control signals,
network latency becomes a factor. To better ascertain the impact of network latency on the centralised FLISR operation, an identical test case was simulated ten times on domestic WiFi, analogous to the
network bandwidth available on 5G networks, 4G LTE, 3G and 2G, with the average time calculated for each network type as presented in \hyperref[table:3]{Table \ref{table:3}}.

\begin{table}
\vspace{0.5cm}
  \caption{\label{table:3}Average centralised FLISR operational speed}
  \centering
  \begin{tabular}{|l|r|}
  \hline
     \header  Network type & Average speed (ms)  \\
     \row Domestic WiFi (5G analogue) & 212.8  \\
     \row 4G LTE & 260.4  \\
     \row 3G & 267.1  \\
     \row 2G & 1049.6 \\
  \hline
  \end{tabular}
\end{table}

\noindent From the findings above, the average operational speed even when using low bandwidth networks such as 2G is quite high, this is significant as when applied to the field,
the most commonly available networks available are likely to be 3G and 2G, particularly in rural areas. When compared to the distributed FLISR algorithm however, network latency is not a
consideration, for the tests conducted the average speed of operation is at or near real-time at two milliseconds.

\begin{itemize}
  \item \textbf{R3 - Given a diverse set of network topologies, will the FLISR solution output valid results for a variety of grid configurations?}
\end{itemize}

\noindent As the findings gathered from the results above have shown, the FLISR Edge solution produces control signals for each of the test cases simulated on the three trial distribution networks. As
stated previously, the effectiveness of the solution depends on the network configuration and the level of simulated damage to the grid, for simpler networks such as Kerry and Mullingar,
the effectiveness of the solution can be more limited compared to a radial network such as Waterford in terms of reducing the number of affected customers and CML accumulated.
However, another aspect of quickly isolating faulted networks is one of safety, where downed lines which are still live pose a safety threat to citizens and therefore the immediate operation
enabled by the FLISR Edge solution would assist in alleviating this danger. A more comprehensive collection of results for both algorithms can be found in \hyperref[sec:appendix_c]{Appendix C}.


% subsection findings (end)
